\documentclass[11pt,a4paper]{article}

% --- 核心宏包与中文支持 ---
\usepackage[utf8]{inputenc}
\usepackage{xeCJK}
\usepackage{amsmath, amssymb, amsfonts}
\usepackage{listings}
\usepackage{xcolor}
\usepackage{geometry}
\usepackage{tcolorbox}
\usepackage{booktabs}
\usepackage{tabularx}
\usepackage{setspace}
\usepackage{enumitem}
\usepackage[hidelinks]{hyperref}
\usepackage{bookmark}
\usepackage{fancyhdr}

% --- PDF 书签与元数据 ---
\hypersetup{
    colorlinks=false,
    pdfborder={0 0 0},
    bookmarks=true,
    bookmarksnumbered=true,
    bookmarksopen=true,
    bookmarksopenlevel=2,
    pdfcreator={XeLaTeX},
    pdftitle={05 - Object Oriented Programming},
    pdfauthor={黄子扬},
}

% --- 页面美化设置 ---
\geometry{top=0.7in,bottom=0.8in,left=0.6in,right=0.6in}
\onehalfspacing
\tcbuselibrary{skins,breakable}
\setlength{\parindent}{0pt}
\setlength{\parskip}{0.6em}

% --- 代码样式配置 ---
\definecolor{mainblue}{RGB}{0,51,102}
\definecolor{examred}{RGB}{180,0,0}
\definecolor{codebg}{RGB}{248,248,248}
\definecolor{commentgreen}{RGB}{0,128,0}

\lstset{
    backgroundcolor=\color{codebg},
    basicstyle=\ttfamily\footnotesize,
    keywordstyle=\color{blue}\bfseries,
    commentstyle=\color{commentgreen},
    stringstyle=\color{orange},
    breaklines=true,
    showstringspaces=false,
    frame=leftline,
    xleftmargin=2em,
    numbers=left,
    numbersep=5pt,
    numberstyle=\tiny\color{gray},
    language=Python,
    morekeywords={None, True, False, class, def, self, super, pass, property, return, if, for, in, del}
}

% --- 自定义教学框 ---
\newtcolorbox{concept}[1]{colback=blue!5, colframe=mainblue, fonttitle=\bfseries, title={{【Concept / 核心概念】#1}}, arc=3pt, breakable}
\newtcolorbox{warning}[1]{colback=red!5, colframe=examred, fonttitle=\bfseries, title={{【Warning / 避坑指南】#1}}, arc=3pt, breakable}
\newtcolorbox{examplebox}[1]{colback=green!2, colframe=green!40!black, fonttitle=\bfseries, title={{【Example / 实操案例】#1}}, arc=2pt, breakable}

% --- 页眉页脚 ---
\pagestyle{fancy}
\fancyhf{}
\fancyhead[L]{\small\bfseries 05 - Object Oriented Programming / 面向对象编程}
\fancyhead[R]{\small\thepage}
\renewcommand{\headrulewidth}{0.5pt}

% --- 文档信息 ---
\title{\Huge \bfseries 05 - Object Oriented Programming}
\author{\Large \textbf{哈尔滨工业大学 \quad 黄子扬}}
\date{2026年6月8日}

\begin{document}

\maketitle
\tableofcontents
\newpage

% ======================================================================
\section{Why OOP? / 为什么需要面向对象？}
% ======================================================================

\begin{concept}{From Lists to Classes / 从列表到类}
    OOP (Object-Oriented Programming) is a method of structuring a program by bundling related properties and behaviors into individual \textbf{objects}. \\
    面向对象编程是一种将相关属性和行为打包成"对象"的程序设计方法。

    \textbf{The Problem with Lists / 用列表存储数据的困境}：

    Imagine tracking employees in an organization. You could represent each employee as a list:
\begin{lstlisting}
Eric = ["Eric", 34, "CEO", 2016]
Alice = ["Alice", "Science Officer", 2017]
\end{lstlisting}
    There are issues with this approach:
    \begin{enumerate}[itemsep=0.3em]
        \item You need to \textbf{remember} which index means what (e.g., index 0 = name, index 1 = age).
        \item What if not every employee has the same number of elements? (Age is missing for Alice!)
        \item What if you want to define \textbf{behaviors} (methods) for your employees?
    \end{enumerate}

    \textbf{Classes} allow you to create \textbf{User-defined Data Structures} that group data (attributes) and behaviors (methods) together. Classes define a type.
\end{concept}

% ======================================================================
\section{Creating Classes / 创建类}
% ======================================================================

\subsection{The Blueprint / 蓝图定义}

\begin{concept}{The class Keyword / class 关键字}
    All class definitions start with the \texttt{class} keyword, followed by the name of the class and a colon (\texttt{:}).
    \begin{itemize}[itemsep=0.3em]
        \item \textbf{Naming / 命名}：Use \textbf{CamelCase} (e.g., \texttt{Employee}), starting with a capital letter. This is different from functions/variables which use \texttt{snake\_case}.
        \item \textbf{\texttt{pass}}：A placeholder used as a body where code will eventually go. Allows you to run the class without throwing an error.
    \end{itemize}
\end{concept}

\begin{examplebox}{Minimal Class / 最简类定义}
\begin{lstlisting}
class Employee:
    pass                # Placeholder — does nothing, but avoids SyntaxError
\end{lstlisting}
\end{examplebox}

\subsection{The \texttt{\_\_init\_\_()} Method / 初始化方法}

\begin{concept}{The Initializer / 初始化器}
    The \texttt{\_\_init\_\_()} method is a special method that runs \textbf{every time} a new object is created. It tells Python what the \textbf{initial state} (initial values of the object's properties) should be. \\
    \texttt{\_\_init\_\_()} 在每次创建新对象时自动运行，设置对象的初始值。
\end{concept}

\begin{warning}{What is self? / self 到底是什么？}
    The first argument of \texttt{\_\_init\_\_()} is \textbf{always} \texttt{self}. It is a reference to the \textbf{specific instance} being created. Think of it as a \textbf{placeholder} while building the blueprint — it refers to the instance that will be created later. \\
    \texttt{self} 指代"当前的这个实例"。它是构建蓝图时的占位符——虽然实例还没创建，但 \texttt{self} 代表了将来会被创建的实例对象。
\end{warning}

% ======================================================================
\section{Attributes / 属性}
% ======================================================================

\begin{concept}{Instance vs Class Attributes / 实例属性 vs 类属性}
    \begin{itemize}[itemsep=0.8em]
        \item \textbf{Instance Attributes / 实例属性}：Belong to a \textbf{specific} object. Defined inside \texttt{\_\_init\_\_()} with \texttt{self.xxx}. Each instance has its own copy. \\
        例：\texttt{self.name} — 每个员工有不同的名字。
        \item \textbf{Class Attributes / 类属性}：Same for \textbf{ALL} instances of a class. Defined directly inside the class body (outside \texttt{\_\_init\_\_()}). All instances share the same value. \\
        例：\texttt{company\_name = "Emlyon"} — 所有员工属于同一公司。
    \end{itemize}
\end{concept}

\begin{examplebox}{Class Definition with Attributes / 完整类定义}
\begin{lstlisting}
class Employee:
    # Class Attribute / 类属性 — shared by all instances
    company_name = "Emlyon"

    def __init__(self, name, age):
        # Instance Attributes / 实例属性 — unique to each instance
        self.name = name
        self.age = age
\end{lstlisting}
\end{examplebox}

% ======================================================================
\section{Instantiating \& Methods / 实例化与方法}
% ======================================================================

\subsection{Instantiating Objects / 实例化对象}

\begin{concept}{Blueprint → Object / 蓝图 → 对象}
    While the \textbf{class} is the blueprint, an \textbf{instance} is an object built from a class that contains real data. \\
    类是蓝图，实例是从类构建出来的包含真实数据的对象。

    To instantiate an object, type the name of the class (CamelCase) followed by parentheses containing any values that must be passed to \texttt{\_\_init\_\_()}.
\end{concept}

\begin{examplebox}{Creating Objects / 创建对象}
\begin{lstlisting}
emp1 = Employee("David", 23)
emp2 = Employee("Alice", 35)

print(emp1)                     # <__main__.Employee object at 0x00aeff70>
# The hex number is the memory address — different on each computer!

print(emp1 == emp2)             # False (different objects in memory)
print(type(emp1))               # <class '__main__.Employee'>

# Access attributes with dot notation / 用点号访问属性
print(emp1.name)                # 'David'
print(emp2.age)                 # 35

# Class attributes are accessed the same way
print(emp1.company_name)        # 'Emlyon'
print(emp2.company_name)        # 'Emlyon' (same for all)
\end{lstlisting}
\end{examplebox}

\begin{warning}{Objects Are Mutable by Default / 对象默认是可变的}
    \begin{lstlisting}
emp1.age = 36                   # Dynamically modify instance attribute
print(emp1.age)                 # 36 — changed!
    \end{lstlisting}
    Custom objects are \textbf{mutable} by default — their attributes can be altered dynamically. This is an important takeaway.
\end{warning}

\subsection{Instance Methods / 实例方法}

\begin{concept}{Behaviors / 行为}
    Classes also have special functions, called \textbf{methods}, that define behaviors and actions that an object can perform with its data. \\
    类内部的函数定义了对象能做的动作。实例方法的第一个参数永远是 \texttt{self}。
\end{concept}

\begin{examplebox}{Instance Methods Example / 实例方法案例}
\begin{lstlisting}
class Employee:
    company_name = "Emlyon"

    def __init__(self, name, age):
        self.name = name
        self.age = age

    # Instance method / 实例方法
    def description(self):
        return f"{self.name} is {self.age} years old"

    # Another instance method / 另一个实例方法
    def update_age(self, new_age):
        self.age = new_age

# Usage
emp1 = Employee("David", 23)
print(emp1.description())       # 'David is 23 years old'
emp1.update_age(24)
print(emp1.description())       # 'David is 24 years old'
\end{lstlisting}
\end{examplebox}

\subsection{Dunder Methods / 魔术方法}

\begin{concept}{The \_\_str\_\_() Method / \_\_str\_\_() 方法}
    By default, \texttt{print(emp1)} outputs something ugly like \texttt{<\_\_main\_\_.Employee object at 0x00aeff70>}. You can change what gets printed by defining a special instance method called \texttt{\_\_str\_\_()}. \\
    定义 \texttt{\_\_str\_\_()} 方法可以改变打印对象时的显示内容。
\end{concept}

\begin{examplebox}{Customizing print() Output / 自定义 print() 输出}
\begin{lstlisting}
class Employee:
    company_name = "Emlyon"

    def __init__(self, name, age):
        self.name = name
        self.age = age

    def __str__(self):
        return f"{self.name} is {self.age} years old"

emp1 = Employee("David", 23)
print(emp1)                     # 'David is 23 years old' — nice!
\end{lstlisting}
\end{examplebox}

\begin{concept}{What Are Dunder Methods? / 什么是魔术方法？}
    Methods like \texttt{\_\_str\_\_()} are called \textbf{dunder methods} because they begin and end with \textbf{d}ouble \textbf{under}scores. There are many dunder methods available that allow your classes to work well with other Python language features (e.g., \texttt{\_\_init\_\_}, \texttt{\_\_str\_\_}, \texttt{\_\_repr\_\_}, \texttt{\_\_eq\_\_}, \texttt{\_\_len\_\_}).
\end{concept}

% ======================================================================
\section{Inheritance / 继承}
% ======================================================================

\begin{concept}{Parent vs Child / 父类与子类}
    \textbf{Inheritance} is the process by which one class takes on the attributes and methods of another. \\
    继承允许一个类获取另一个类的所有属性和方法。

    \begin{itemize}[itemsep=0.5em]
        \item \textbf{Parent Class / 父类}：The class being inherited from.
        \item \textbf{Child Class / 子类}：The class that inherits. Child classes can \textbf{override} (redefine) and \textbf{extend} (add to) the parent's attributes and methods.
        \item To create a child class, put the parent class name in parentheses after the child class name.
    \end{itemize}
\end{concept}

\begin{concept}{The \texttt{super()} Function / super() 函数}
    \texttt{super()} calls the \textbf{parent} class's methods, especially useful for calling \texttt{\_\_init\_\_()} in the child class to initialize inherited attributes. \\
    \texttt{super()} 用于调用父类的方法，最常见的用法是在子类的 \texttt{\_\_init\_\_()} 中调用父类的 \texttt{\_\_init\_\_()}。
\end{concept}

\begin{examplebox}{Full Inheritance Example / 继承完整案例}
\begin{lstlisting}
# Parent class / 父类
class Employee:
    company_name = "Emlyon"

    def __init__(self, name, age):
        self.name = name
        self.age = age

    def __str__(self):
        return f"{self.name} is {self.age} years old"

# Child class 1: HR Employee / 子类 1
class HREmployee(Employee):
    def __init__(self, name, age, wh):
        super().__init__(name, age)     # Call parent's __init__
        self.working_hours = wh         # New attribute / 新属性

    def update_working_hours(self, new_hour):   # New method / 新方法
        self.working_hours = new_hour

    def __str__(self):                          # Override __str__ / 重写
        return f"{self.name} is {self.age} years old and works in HR and has {self.working_hours} working hours!"

# Child class 2: Management Employee / 子类 2
class ManagementEmployee(Employee):
    def __init__(self, name, age):
        super().__init__(name, age)
        self.emp_list = []              # New attribute / 新属性

    def add_employee(self, emp_name):   # New method / 新方法
        self.emp_list.append(emp_name)

    def give_emp_list(self):            # New method / 新方法
        for i in range(len(self.emp_list)):
            print(f"{i} : {self.emp_list[i]}")
\end{lstlisting}
\end{examplebox}

\begin{examplebox}{Using the Child Classes / 使用子类}
\begin{lstlisting}
# HR Employee / HR 员工
hr_emp_1 = HREmployee("David", 23, 8)
print(hr_emp_1)                     # Uses overridden __str__
hr_emp_1.update_working_hours(10)   # Child's own method
print(hr_emp_1.working_hours)       # 10

# Management Employee / 管理员工
management_emp_1 = ManagementEmployee("Alice", 35)
print(management_emp_1)             # Uses parent's __str__
management_emp_1.add_employee("John")
management_emp_1.add_employee("Maria")
management_emp_1.give_emp_list()
# Output:
# 0 : John
# 1 : Maria
\end{lstlisting}
\end{examplebox}

% ======================================================================
\section{Encapsulation \& Properties / 封装与属性装饰器}
% ======================================================================

\begin{concept}{Read-Only Attributes / 只读属性}
    Sometimes it's a bad idea if someone can change an attribute from outside the class, but you still want to \textbf{see} the value. This is done using a \textbf{getter-function} with the \textbf{\texttt{@property}} decorator. \\
    使用 \texttt{@property} 装饰器可以将方法伪装成属性，实现只读访问，防止外部非法篡改。
\end{concept}

\begin{examplebox}{@property Example / @property 案例}
\begin{lstlisting}
class MyClass:
    def __init__(self):
        self._hidden = 42           # Prefix '_' suggests internal use

    @property
    def hidden(self):               # Getter — can access like an attribute
        return self._hidden

    def get_hidden(self):           # Ordinary getter — requires parentheses
        return self._hidden

my_object = MyClass()

# @property: access like an attribute (no parentheses!)
print(my_object.hidden)             # 42

# Direct assignment to @property is BLOCKED
# my_object.hidden = 0              # AttributeError! Read-only!

# BUT: _hidden can still be accessed directly if you really want
my_object._hidden = 0               # Works, but you shouldn't do it
print(my_object._hidden)            # 0

# Ordinary getter requires parentheses
print(my_object.get_hidden())       # 0
\end{lstlisting}
\end{examplebox}

\begin{warning}{The \_ Prefix Convention / \_ 前缀约定}
    The underscore prefix (\texttt{\_hidden}) is a \textbf{convention} that tells other programmers "this attribute is internal — don't access it directly." However, it does \textbf{NOT} actually prevent access. Python trusts you to follow the convention.
\end{warning}

% ======================================================================
\section{Quick Reference / 速查表}
% ======================================================================

\begin{table}[h]
\centering
\small
\renewcommand{\arraystretch}{1.3}
\begin{tabular}{@{}ll@{}}
\toprule
\textbf{Concept / 概念} & \textbf{Syntax / 语法} \\ \midrule
定义类 & \texttt{class MyClass:} \\
初始化方法 & \texttt{def \_\_init\_\_(self, param):} \\
实例属性 & \texttt{self.attr = value} \\
类属性 & \texttt{class\_attr = value}（在 \texttt{\_\_init\_\_} 外） \\
实例化 & \texttt{obj = MyClass(arg)} \\
字符串表示 & \texttt{def \_\_str\_\_(self): return "..."} \\
继承 & \texttt{class Child(Parent):} \\
调用父类 & \texttt{super().method()} \\
只读属性 & \texttt{@property} 装饰 getter 方法 \\
内部属性约定 & \texttt{self.\_internal} \\ \bottomrule
\end{tabular}
\end{table}

\end{document}
